% MAREA validation v3: the Cantabrian estuaries against the PNOA LiDAR. Its tables, figures and macros are written
% by build_v3.py, which compiles it.

\documentclass[10pt,a4paper]{article}
\usepackage[T1]{fontenc}
\usepackage[utf8]{inputenc}
\usepackage[spanish,es-noshorthands,es-nodecimaldot,es-tabla]{babel}
\usepackage{lmodern}
\usepackage[a4paper,top=2.2cm,bottom=2.4cm,left=2.0cm,right=2.0cm,headsep=0.6cm]{geometry}
\usepackage{amsmath}
\usepackage{graphicx,booktabs,array,multirow,enumitem,float}
\usepackage{xcolor}
\usepackage{caption}
\usepackage{microtype}
\usepackage{fancyhdr}
\usepackage[hidelinks]{hyperref}

\definecolor{accent}{HTML}{B5523B}
\definecolor{ink}{HTML}{2B2B2B}
\definecolor{secondary}{HTML}{6B6B6B}
\color{ink}

\newcommand{\valdef}[2]{\expandafter\def\csname val@#1\endcsname{#2}}
\DeclareRobustCommand{\val}[1]{\ifcsname val@#1\endcsname\csname val@#1\endcsname\else
  \PackageError{report}{Missing value #1}{build_v3.py did not write it}\fi}
\input{tables/macros.tex}

\captionsetup{font=small,labelfont=bf,labelsep=period,justification=raggedright,singlelinecheck=false}
\captionsetup[table]{position=top,skip=3pt}
\setlength{\parindent}{0pt}
\setlength{\parskip}{0.5ex plus 0.2ex}
\pagestyle{fancy}
\fancyhf{}
\renewcommand{\headrulewidth}{0pt}
\fancyhead[L]{\footnotesize\color{secondary}MAREA · validación v3 · estuarios cantábricos}
\fancyhead[R]{\footnotesize\color{secondary}borrador interno · ningún resultado es definitivo}
\fancyfoot[C]{\footnotesize\color{secondary}\thepage}

\begin{document}

{\Large\bfseries Validación v3: estuarios cantábricos frente al LiDAR del PNOA}\\[2pt]
{\small\color{secondary}Compilado el \val{build:date}. Verdades con resultados: \val{truths}.}

\section*{Qué se valida}

\textbf{Verdad.} Los modelos del terreno del PNOA-LiDAR que indicó F.~García-González: MDT02 de la 2.ª cobertura
(Villaviciosa y Eo 2020, Foz 2016, Mundaka 2017) y MDT50CM de la 3.ª (Mundaka 2024); el año es el que da el CNIG para
cada fichero. Se enmascara el agua en el vuelo (parches planos de al menos 200\,m²) y se promedia a la malla de
10\,m de Sentinel-2. Fecha de verdad provisional: 15 de julio del año del vuelo.

\textbf{Ventanas.} Para cada $n$ = 20, 40, 60, 80, 100 y 120, las $n$ escenas cribadas más cercanas a la fecha de
verdad, simétricas en el tiempo. Si el cubo no admite la ventana centrada en la fecha de verdad (el archivo de
Sentinel-2 empieza en julio de 2015), su centro se desplaza los días mínimos para que quepa, con la fecha de verdad
dentro (Foz); el desplazamiento figura en la tabla de ventanas de cada verdad. Una ventana que no cabe de ningún modo
no es alcanzable.

\textbf{Píxeles.} Los de la extracción de producción con verdad dentro del rango del producto sin retardo
($\pm 0{,}25$\,m), y solo aquellos en los que ninguna clase domina la ventana: fracción de miradas mojadas entre
0,10 y 0,90 ($q = 0{,}10$; la sensibilidad a $q$ va en la figura~\ref{fig:sens}).

\textbf{Métodos.} Los diez de las tablas $n$, MAREA-B, MAREA-BL y MAREA-S (un único ajuste, retardo suavizado).

\begin{figure}[H]\centering
\includegraphics[width=\textwidth]{figures/curves_v3.pdf}
\caption{Error frente a $n$ en cada verdad, píxeles propios de cada método ($q = 0{,}10$). MAREA y sus versiones en
color; el resto en gris.}
\end{figure}

\IfFileExists{figures/sensitivity_v3.pdf}{%
\begin{figure}[H]\centering
\includegraphics[width=\textwidth]{figures/sensitivity_v3.pdf}
\caption{Sensibilidad a la máscara de clase dominante: RMSE en los píxeles comunes frente a $q$ (se puntúan los
píxeles con fracción mojada entre $q$ y $1-q$), en la ventana mayor de cada verdad; la línea punteada es el valor
adoptado, $q = 0{,}10$. Colores y trazos como en la figura~1.}\label{fig:sens}
\end{figure}}{}

\clearpage
\section*{Tablas por verdad}
\input{tables/body.tex}

\section*{Mapas de error y descomposición}
Por verdad, en $n = 100$ y en la ventana menor calculada: el error $p - t$ de cada método en cada píxel (escala
divergente centrada en 0, la misma para todos los métodos y ventanas de una verdad) y la descomposición del error
cuadrático medio, $\mathrm{MSE} = \overline{d}^{\,2} + ((1-\beta)\,s_t)^2 + s_e^2$: desplazamiento, compresión del
relieve y dispersión.

\input{tables/maps.tex}

\clearpage
\section*{Mapas de cota sin verdad: Santander y Santoña}
Su LiDAR se voló en pleamar, así que no hay puntuación: solo la cota de cada método, en la ventana de las escenas
cribadas más cercanas al 15 de julio de 2024 ($n = 100$ en Santoña; $n = 40$ en Santander, cuyo cubo de 2023--2025
solo tiene 51 escenas cribadas).

\input{tables/elevation_maps.tex}

\section*{Avisos sobre la verdad}
\begin{itemize}[leftmargin=1.2em,itemsep=2pt]
\item \textbf{Agua en el vuelo.} En el intermareal bajo el LiDAR mide la superficie del agua donde la había durante el
vuelo, no el fondo. Se aprecia en los valores repetidos de la verdad baja (Villaviciosa, Foz), en el solape de las
hojas 0010-3 y 0010-4 del Eo (difieren en el canal y coinciden en lo alto) y, sobre todo, en Mundaka: la 3.ª cobertura
queda casi un metro por encima de la 2.ª en lo más bajo y coincide con ella en el terreno estable. La máscara de nivel
de vuelo del diseño lo resolvería, pero necesita la hora del vuelo (las nubes de puntos LAZ la llevan).
\item \textbf{Píxeles que no sigue la marea.} En Villaviciosa hay una zona (charcas o marisma con compuertas) donde
MAREA-B, MAREA-S y SE dan cotas más de un metro por debajo del LiDAR; DEA y MAREA V0 se abstienen allí. Por eso el
orden cambia entre los píxeles propios de cada método y los comunes.
\end{itemize}

\end{document}
